Linear models such as DLinear were reported to match or beat Transformer forecasters on common long-horizon benchmarks, but real datasets do not let one vary the data property that might favour nonlinear models. We compare, under one protocol, a seasonal-naive forecaster, a DLinear-style decomposition-plus-linear model, a small patch transformer (PatchTST-style) and a small patch GRU, all reimplemented on CPU, at horizons 24, 96 and 192. The data are synthetic series with controlled seasonality, trend, noise and Markov regime switching, plus ETTh1; three seeds, no hyper-parameter tuning. A win is defined in advance as at least 5\% lower seed-mean MSE and lower MSE in all three seeds. The answer depends on the training budget. At the planned 400 training steps neither nonlinear model wins on any of the six tasks at any horizon; the closest case is the patch transformer on the regime-switching task (4.8\% better at horizon 96). A training-budget sweep, added after an independent audit, changes this on the regime-switching task only: at 1000 and 2000 steps both nonlinear models win at horizon 24 and the GRU wins at horizon 96 (10.6\% better at 2000 steps), while DLinear's MSE barely moves with the budget; nobody wins at horizon 192. On the four tasks without regime changes (tested at horizon 96) and on ETTh1 (all horizons) no nonlinear model wins at 2000 steps either, and on ETTh1 their MSE is 2.9\% to 9.3\% above DLinear's at horizons 96 and 192. At 2000 steps the GRU passes the rule at mean regime lengths of 100 and 500 steps but not at 250, 1000 or 4000; its margin is not monotone in the switching rate. Removing the decomposition changes DLinear by less than 2\% (400 steps). The evidence is small-scale: three seeds, at most 2000 steps, untuned models and synthetic generators of our own design; the noise sweep and the design ablation were run at 400 steps only.
